% !TeX root = ../PIFB2.tex
\section{Conclusion}\label{sec:global_conclusion}

A participatory theory needs a disciplined account of partial access. Local statistical sections over a continuous base, connected through selective matrix-valued channels, supply such an account. One normalized population law yields a collective evidence bound and source-selection attention. One joint statistical bundle accommodates state beliefs and model uncertainty. Exact coarse-graining produces parent laws, induced interactions and conditional cross-scale operators; predictive updating supplies later priors without a same-time shadow postulate.

The framework also gives concrete limits. Coarse-graining does not select its own retained variables, a raw VFE decrease does not certify formation, and slow memory does not establish recursive coalescence. Fisher length is an information clock, not yet physical time. Hermitian determinant geometry carries the physical Lorentz group's kinematics under explicit assumptions, but its realization as spacetime remains to be derived. Persistent dissipative structures and memory can emerge without a kinetic postulate, while physical mass requires further dynamical structure.

These distinctions preserve Wheeler's ambition while making progress assessable. The next results should identify which collective variables the actual dynamics retain, how their laws compose, and what physical observations those laws predict.

